\section{Limiting cases and counterexamples}
\label{sec:limits}

Every compact formula in this paper has a nearby case in which its geometry
collapses. The cases below are tests, not footnotes.

\begin{description}
 \item[\(R_s\to0\).] The PSM voltage cross term vanishes and its axes align
 with \(dq\). The exact center becomes the characteristic point at every
 nonzero speed. In the EESM, the voltage null direction becomes
 \([-L_{\mathrm{exc}}/L_d,0,1]^{\mathsf T}\). Loss whitening becomes singular
 only if the same idealization is also applied to stator copper loss; therefore
 a zero-resistance voltage approximation must not automatically replace the
 positive thermal resistance.

 \item[\(\omega_e\to0\).] With \(R_s>0\), the PSM voltage ellipse tends to a
 resistance circle centered at the origin. For the EESM, excitation current is
 invisible to stator voltage at exactly zero speed, while copper loss still
 bounds it. If also \(R_s=0\), the voltage map is zero and the voltage
 constraint imposes no current restriction.

 \item[\(|\omega_e|\to\infty\).] The PSM center tends to
 \((-\psi_{\mathrm{pm}}/L_d,0)\), its axes scale approximately as
 \(1/|\omega_e|\), and the sign of its small \(q\)-center component reverses
 with speed. The EESM voltage cylinder narrows in its two visible directions
 but remains open along its null line.

 \item[\(L_d=L_q\).] The PSM voltage boundary is a translated circle and the
 rotation angle is undefined, not zero by decree. Reluctance torque vanishes,
 but EESM excitation torque remains. In active-flux coordinates the effective
 coefficient vector simply loses its \(d\)-current component.

 \item[\(\psi_{\mathrm{pm}}=0\).] The PSM voltage ellipse returns to the origin.
 Torque is purely reluctance torque and vanishes altogether in the additional
 nonsalient limit. The EESM parent remains meaningful because controllable
 excitation can replace the missing PM flux.

 \item[Fixed \(i_{\mathrm{exc}}\).] A homogeneous 3D EESM problem becomes an
 affine 2D PSM-like slice. Rotor copper loss is a constant slice offset; an
 actual PSM omits it.

 \item[\(P_{\mathrm{Cu},\max}\to\infty\).] The linear EESM maximum-torque
 problem under voltage alone is generically unbounded. Following the voltage
 null line leaves voltage unchanged while torque can grow. A finite numerical
 box would manufacture a false optimum.

 \item[\(U_{\max}\to\infty\).] Voltage disappears from allocation. The exact
 solution is \(z_0=0\), \(|z_\psi|=|z_q|\), and
 \(P_{\min}=2|M|/\kappa\). Maximum torque under a loss budget follows the
 corresponding dominant torque-per-loss eigenvector.

 \item[Vanishing saliency.] With \(\Delta L=0\), direct-axis stator current is
 torque-neutral in the linear EESM torque law; field excitation alone creates
 active flux. With vanishing excitation coupling instead, field current is
 torque-neutral and the machine reduces to reluctance torque allocation.

 \item[Degenerate eigenvalues.] Principal directions become nonunique.
 Algorithms must compare invariant subspaces or residuals rather than raw
 eigenvector angles.

 \item[Positive and negative torque.] The loss and voltage sets are even in
 the full current vector, so \(\vect i\mapsto-\vect i\) does not reverse torque.
 Torque reversal requires changing the sign of one product factor, producing
 the opposite saddle branch. Extra one-sided field constraints can make motor
 and generator capability asymmetric.

 \item[Unreachable request.] The voltage-limited direction set may contain no
 point capable of \(M_{\mathrm{ref}}\), or its minimum loss may exceed the
 thermal ceiling. The controller must saturate on \(M_{\max}\) and return a
 limit status; scaling currents beyond the boundary is not a solution.
\end{description}

These limits are also debugging oracles. Any implementation that reports a
unique ellipse angle in the circle case, a bounded voltage-only EESM torque,
or a homogeneous 2D PSM model with nonzero fixed PM flux has contradicted the
geometry before numerical tolerances are considered.
